\clearpage
\fancyhead[L]{\small BEIMA. DE ANNULO SATURNI.}
\hypertarget{native-659-middle}{}
\begin{center}\small Göttingische gelehrte Anzeigen. Stück 194. Seiten 1937..1938. 5. Dezember 1842.\\[3pt]\rule{.72\linewidth}{.3pt}\end{center}

\emph{De annulo Saturni commentatus est} \textsc{Elto Martini Beima}, \emph{mus. hist. nat. publ. conservator.} Leyden 1842. Bei S. und J. Luchtmans. 252 Seiten im Quartformat, dazu 4 Kupfertafeln.

In dieser Monografie über einen der interessantesten Gegenstände unseres Sonnensystems ist mit vielem Fleiß alles gesammelt, was die Geschichte der Entdeckung des Saturnrings betrifft, und was die Astronomen über seine Lage, Gestalt, Dimensionen, Masse, physikalische Beschaffenheit, Rotation und Ursprung sowie über die Kräfte, die ihn erhalten, und über die Erscheinungen, die er für seinen Hauptplaneten darbietet, bisher beobachtet, berechnet oder vermutet haben. Höhere Ansprüche, als auf den Rang einer mit Fleiß erarbeiteten und aus den Quellen geschöpften Kompilation kann die Schrift nicht erheben, aber dieses Lob kann man ihr zugestehen, wenn man gegenüber einer sehr fehlerhaften und nachlässigen Sprache nachsichtig ist, die stellenweise kaum ohne Heranziehung der Quellen, aus denen übersetzt wurde, verständlich wird.
